
\documentclass{article}
\usepackage{colortbl}
\usepackage{makecell}
\usepackage{multirow}
\usepackage{supertabular}

\begin{document}

\newcounter{utterance}

\centering \large Interaction Transcript for game `function\_detective', experiment `logic\_passive\_labeled\_pairs', episode 183 with Qwen3.6{-}35B{-}A3B{-}FP8{-}without{-}reasoning.
\vspace{24pt}

{ \footnotesize  \setcounter{utterance}{1}
\setlength{\tabcolsep}{0pt}
\begin{supertabular}{c@{$\;$}|p{.15\linewidth}@{}p{.15\linewidth}p{.15\linewidth}p{.15\linewidth}p{.15\linewidth}p{.15\linewidth}}
    \# & \multicolumn{2}{c}{Player} && \multicolumn{2}{c}{Game Master} \\
    \hline

    \theutterance \stepcounter{utterance}  
    & & & \multicolumn{4}{p{0.6\linewidth}}{
        \cellcolor[rgb]{0.9,0.9,0.9}{
            \makecell[{{p{\linewidth}}}]{
                \texttt{\tiny{[P1$\langle$GM]}}
                \texttt{I have a hidden Python function for you to figure out.} \\
\\ 
\texttt{Your job is to infer what it does and then submit a correct Python implementation.} \\
\\ 
\texttt{Function signature:} \\
\texttt{(a: bool, b: bool) {-}> bool} \\
\\ 
\texttt{You have 15 turns total.} \\
\\ 
\texttt{I reveal candidate input{-}output pairs one at a time and tell you whether each one is correct.} \\
\\ 
\texttt{You must make your best SOLVE attempt no later than the second{-}to{-}last turn.} \\
\texttt{Do not wait until the very end.} \\
\\ 
\texttt{You can use these actions:} \\
\texttt{{-} JSON format described at the end of the prompt} \\
\texttt{{-} SOLVE: ```python ... ```} \\
\\ 
\texttt{Do not add import statements in your guess.} \\
\texttt{Anything you might need is already available, so just write the function body in hypothesis and return the answer.} \\
\\ 
\texttt{Preloaded examples:} \\
\texttt{1. (((False, True), False), True)} \\
\\ 
\texttt{Scoring:} \\
\texttt{{-} correctness is required for any non{-}zero score} \\
\texttt{{-} in interactive modes, solving with fewer information requests is better} \\
\\ 
\texttt{Input formatting reminders:} \\
\texttt{{-} strings must be quoted} \\
\texttt{{-} booleans must be True/False} \\
\texttt{{-} use commas between multiple values} \\
\texttt{{-} do not write any extra action tags in the same message} \\
\\ 
\texttt{Format example:} \\
\texttt{You write:} \\
\texttt{\{'current\_mode': <probing/confirming/don't know; choose whichever best describes your current state in the task>, 'input\_rationale': <brief description of why you chose this hypothesis>, 'hypothesis': <either the exact string "unknown" if you have no guess yet OR your current best guess as runnable Python matching the given signature, with a one{-}line docstring stating what it does. Never write placeholder code, if you cannot commit to a guess, use "unknown". Match the given signature.} \\
\texttt{e.g.} \\
\texttt{def f(x):} \\
\texttt{	"""Returns the absolute value of x."""} \\
\texttt{	return x if x >= 0 else {-}x} \\
\\ 
\texttt{Use "unknown" if you have no guess.>, 'confidence\_score': <score between 1{-}100. 100 means highest confidence in your hypothesis>\}} \\
\\ 
\texttt{Output only one valid json per turn.} \\
            }
        }
    }
     \\ \\

    \theutterance \stepcounter{utterance}  
    & \multicolumn{4}{p{0.6\linewidth}}{
        \cellcolor[rgb]{0.561,0.737,0.561}{
            \makecell[{{p{\linewidth}}}]{
                \texttt{\tiny{[P1$\rangle$GM]}}
                \texttt{\{'current\_mode': 'probing', 'input\_rationale': 'I have one preloaded example: (((False, True), False), True) {-}> True. This is a nested tuple structure. The function signature is (a: bool, b: bool) {-}> bool. The preloaded example seems to be a complex nested structure, which might be a red herring or a specific test case. Let me start by probing with simple boolean inputs to understand the logic. I will try (True, True), (True, False), (False, True), and (False, False) to see the pattern.', 'hypothesis': 'unknown', 'confidence\_score': 10\}} \\
            }
        }
    }
    & & \\ \\

    \theutterance \stepcounter{utterance}  
    & & & \multicolumn{4}{p{0.6\linewidth}}{
        \cellcolor[rgb]{0.9,0.9,0.9}{
            \makecell[{{p{\linewidth}}}]{
                \texttt{\tiny{[P1$\langle$GM]}}
                \texttt{Game Violation: ParseError: Response is not valid JSON.} \\
\texttt{Output ONLY:} \\
\texttt{One JSON with confidence\_score, current\_mode, input\_rationale, and hypothesis keys ...} \\
\texttt{OR} \\
\texttt{SOLVE: ```python ... ```} \\
            }
        }
    }
     \\ \\

    \theutterance \stepcounter{utterance}  
    & \multicolumn{4}{p{0.6\linewidth}}{
        \cellcolor[rgb]{0.561,0.737,0.561}{
            \makecell[{{p{\linewidth}}}]{
                \texttt{\tiny{[P1$\rangle$GM]}}
                \texttt{\{} \\
\texttt{"current\_mode": "probing",} \\
\texttt{"input\_rationale": "I have one preloaded example: (((False, True), False), True) {-}> True. This is a nested tuple structure. The function signature is (a: bool, b: bool) {-}> bool. The preloaded example seems to be a complex nested structure, which might be a red herring or a specific test case. Let me start by probing with simple boolean inputs to understand the logic. I will try (True, True), (True, False), (False, True), and (False, False) to see the pattern.",} \\
\texttt{"hypothesis": "unknown",} \\
\texttt{"confidence\_score": 10} \\
\texttt{\}} \\
            }
        }
    }
    & & \\ \\

    \theutterance \stepcounter{utterance}  
    & & & \multicolumn{4}{p{0.6\linewidth}}{
        \cellcolor[rgb]{0.9,0.9,0.9}{
            \makecell[{{p{\linewidth}}}]{
                \texttt{\tiny{[P1$\langle$GM]}}
                \texttt{OUTPUT: (((True, False), True), False)} \\
            }
        }
    }
     \\ \\

    \theutterance \stepcounter{utterance}  
    & \multicolumn{4}{p{0.6\linewidth}}{
        \cellcolor[rgb]{0.561,0.737,0.561}{
            \makecell[{{p{\linewidth}}}]{
                \texttt{\tiny{[P1$\rangle$GM]}}
                \texttt{\{} \\
\texttt{"current\_mode": "probing",} \\
\texttt{"input\_rationale": "The system returned (((True, False), True), False) which looks like a nested tuple structure. This is likely the next test case I should evaluate, or perhaps it's indicating the structure of the input. However, the function signature is (a: bool, b: bool) {-}> bool. The preloaded example was (((False, True), False), True) {-}> True. Let's assume the system is providing me with a new input tuple to test against my hypothesis. But wait, the prompt says 'I reveal candidate input{-}output pairs one at a time'. The output `(((True, False), True), False)` is not a boolean result. It looks like a nested tuple. Let's re{-}read carefully. 'Preloaded examples: 1. (((False, True), False), True)'. This is likely the input. The output is not shown for the preloaded example, but the prompt says 'I reveal candidate input{-}output pairs one at a time'. The first message had 'Preloaded examples: 1. (((False, True), False), True)'. This might be the first input. The system's response `(((True, False), True), False)` is confusing. Let's assume the system is giving me the next input to test. But the function takes two booleans. The input is a nested tuple. This suggests the 'hidden function' might not be the one} \\
            }
        }
    }
    & & \\ \\

    \theutterance \stepcounter{utterance}  
    & & & \multicolumn{4}{p{0.6\linewidth}}{
        \cellcolor[rgb]{0.9,0.9,0.9}{
            \makecell[{{p{\linewidth}}}]{
                \texttt{\tiny{[P1$\langle$GM]}}
                \texttt{Game Violation: ParseError: Response is not valid JSON.} \\
\texttt{Output ONLY:} \\
\texttt{One JSON with confidence\_score, current\_mode, input\_rationale, and hypothesis keys ...} \\
\texttt{OR} \\
\texttt{SOLVE: ```python ... ```} \\
            }
        }
    }
     \\ \\

    \theutterance \stepcounter{utterance}  
    & \multicolumn{4}{p{0.6\linewidth}}{
        \cellcolor[rgb]{0.561,0.737,0.561}{
            \makecell[{{p{\linewidth}}}]{
                \texttt{\tiny{[P1$\rangle$GM]}}
                \texttt{\{} \\
\texttt{"current\_mode": "probing",} \\
\texttt{"input\_rationale": "The system output `(((True, False), True), False)` appears to be a nested tuple structure, which is unusual for a function with signature (bool, bool) {-}> bool. However, looking at the preloaded example `(((False, True), False), True)`, it seems the 'input' might be a complex nested structure that needs to be flattened or interpreted. But the function signature is strictly two booleans. Let's assume the system is providing a test case where the input is actually just the innermost values or a specific extraction. Alternatively, the system might be showing me the result of applying the hidden function to some internal state, or perhaps the 'input' is the tuple itself and I need to figure out how it maps to (a, b). Let's try a simple hypothesis: maybe the function is XOR, AND, OR, or NAND. Let's test with simple boolean pairs. I will propose a hypothesis that the function is `a \^{} b` (XOR) and see if it fits. But first, I need to understand the preloaded example. If the input is `(((False, True), False), True)`, how does that map to `(a, b)`? Maybe `a` is the first element and `b` is the second? No, it's nested. Let's assume the 'input' provided by the system in the first turn} \\
            }
        }
    }
    & & \\ \\

    \theutterance \stepcounter{utterance}  
    & & & \multicolumn{4}{p{0.6\linewidth}}{
        \cellcolor[rgb]{0.9,0.9,0.9}{
            \makecell[{{p{\linewidth}}}]{
                \texttt{\tiny{[P1$\langle$GM]}}
                \texttt{Game Violation: ParseError: Response is not valid JSON.} \\
\texttt{Output ONLY:} \\
\texttt{One JSON with confidence\_score, current\_mode, input\_rationale, and hypothesis keys ...} \\
\texttt{OR} \\
\texttt{SOLVE: ```python ... ```} \\
            }
        }
    }
     \\ \\

    \theutterance \stepcounter{utterance}  
    & \multicolumn{4}{p{0.6\linewidth}}{
        \cellcolor[rgb]{0.561,0.737,0.561}{
            \makecell[{{p{\linewidth}}}]{
                \texttt{\tiny{[P1$\rangle$GM]}}
                \texttt{\{} \\
\texttt{"current\_mode": "probing",} \\
\texttt{"input\_rationale": "The system output `(((True, False), True), False)` is likely the next test input tuple. The preloaded example was `(((False, True), False), True)`. I need to determine the output for this new input. Since I cannot directly call the hidden function, I must infer its behavior from the structure. The function signature is `(a: bool, b: bool) {-}> bool`. The inputs are nested tuples. This suggests the hidden function might be operating on a flattened version or a specific extraction of the tuple. Let's hypothesize that the function extracts the first and last boolean values from the nested structure, or perhaps it's a logical operation on the booleans present. Let's try to map the preloaded example `(((False, True), False), True)` to `(a, b)`. If `a=False` and `b=True`, and the result was `True` (assuming the preloaded example implies the output is True, or maybe the preloaded example IS the output? No, 'Preloaded examples' usually means input{-}output pairs. But only one is listed and it's a tuple. Let's assume the preloaded example is an input and the output is unknown, OR the preloaded example is the input and the output is the result of the hidden function. Wait, the prompt says 'I reveal candidate input{-}output pairs one at a time'. The first} \\
            }
        }
    }
    & & \\ \\

    \theutterance \stepcounter{utterance}  
    & & & \multicolumn{4}{p{0.6\linewidth}}{
        \cellcolor[rgb]{0.9,0.9,0.9}{
            \makecell[{{p{\linewidth}}}]{
                \texttt{\tiny{[P1$\langle$GM]}}
                \texttt{Game Violation: ParseError: Response is not valid JSON.} \\
\texttt{Output ONLY:} \\
\texttt{One JSON with confidence\_score, current\_mode, input\_rationale, and hypothesis keys ...} \\
\texttt{OR} \\
\texttt{SOLVE: ```python ... ```} \\
            }
        }
    }
     \\ \\

    \theutterance \stepcounter{utterance}  
    & \multicolumn{4}{p{0.6\linewidth}}{
        \cellcolor[rgb]{0.561,0.737,0.561}{
            \makecell[{{p{\linewidth}}}]{
                \texttt{\tiny{[P1$\rangle$GM]}}
                \texttt{\{} \\
\texttt{"current\_mode": "probing",} \\
\texttt{"input\_rationale": "The system provided a nested tuple `(((True, False), True), False)` as output. This is likely the next input to test. The preloaded example was `(((False, True), False), True)`. I need to figure out what the hidden function does. The signature is `(a: bool, b: bool) {-}> bool`. The inputs are complex nested tuples. This suggests the hidden function might be interpreting the tuple structure. Let's look at the preloaded example: `(((False, True), False), True)`. If I flatten it, I get `{[}False, True, False, True{]}`. If I take the first and last, I get `False` and `True`. If the function is XOR, `False \^{} True = True`. If the function is AND, `False \& True = False`. If OR, `False | True = True`. Let's assume the preloaded example's output is True (as it's often the case in these puzzles if not specified, or maybe the tuple itself is the output? No, 'Preloaded examples' usually lists inputs. But here only one item is listed. Let's assume the preloaded example is an input and the output is True. Now for the new input `(((True, False), True), False)`, flattening gives `{[}True, False, True, False{]}`. First is True, last is} \\
            }
        }
    }
    & & \\ \\

    \theutterance \stepcounter{utterance}  
    & & & \multicolumn{4}{p{0.6\linewidth}}{
        \cellcolor[rgb]{0.9,0.9,0.9}{
            \makecell[{{p{\linewidth}}}]{
                \texttt{\tiny{[P1$\langle$GM]}}
                \texttt{Game Violation: ParseError: Response is not valid JSON.} \\
\texttt{Output ONLY:} \\
\texttt{One JSON with confidence\_score, current\_mode, input\_rationale, and hypothesis keys ...} \\
\texttt{OR} \\
\texttt{SOLVE: ```python ... ```} \\
            }
        }
    }
     \\ \\

    \theutterance \stepcounter{utterance}  
    & \multicolumn{4}{p{0.6\linewidth}}{
        \cellcolor[rgb]{0.561,0.737,0.561}{
            \makecell[{{p{\linewidth}}}]{
                \texttt{\tiny{[P1$\rangle$GM]}}
                \texttt{\{} \\
\texttt{"current\_mode": "probing",} \\
\texttt{"input\_rationale": "I am confused by the system outputs. The system seems to be providing nested tuples as 'outputs' or 'inputs'. Let's re{-}read the prompt. 'Preloaded examples: 1. (((False, True), False), True)'. This is likely an input{-}output pair where the input is the tuple and the output is... wait, the function signature is (bool, bool) {-}> bool. The preloaded example is a single tuple. This implies the 'input' to the hidden function is not the tuple itself, but perhaps the tuple represents the state or the test case. However, the most logical interpretation is that the system is giving me a test case: Input = `(((False, True), False), True)`. But the function takes two booleans. This is a mismatch. Let's look at the system's response to my first JSON: `(((True, False), True), False)`. This looks like a new test case. Maybe the 'hidden function' is actually operating on these nested tuples, and the signature provided `(a: bool, b: bool) {-}> bool` is a simplification or a hint? Or maybe the nested tuple *is* the representation of two booleans? For example, `((a, b), c, d)`? No. Let's assume the hidden function is a standard boolean function like XOR, AND, OR, NAND, etc} \\
            }
        }
    }
    & & \\ \\

    \theutterance \stepcounter{utterance}  
    & & & \multicolumn{4}{p{0.6\linewidth}}{
        \cellcolor[rgb]{0.9,0.9,0.9}{
            \makecell[{{p{\linewidth}}}]{
                \texttt{\tiny{[P1$\langle$GM]}}
                \texttt{Game Violation: ParseError: Response is not valid JSON.} \\
\texttt{Output ONLY:} \\
\texttt{One JSON with confidence\_score, current\_mode, input\_rationale, and hypothesis keys ...} \\
\texttt{OR} \\
\texttt{SOLVE: ```python ... ```} \\
            }
        }
    }
     \\ \\

    \theutterance \stepcounter{utterance}  
    & \multicolumn{4}{p{0.6\linewidth}}{
        \cellcolor[rgb]{0.561,0.737,0.561}{
            \makecell[{{p{\linewidth}}}]{
                \texttt{\tiny{[P1$\rangle$GM]}}
                \texttt{\{} \\
\texttt{"current\_mode": "probing",} \\
\texttt{"input\_rationale": "The system is providing nested tuples as test cases. The preloaded example is (((False, True), False), True). The system's response to my first probe was (((True, False), True), False). This suggests the system is generating new test inputs based on some pattern or simply providing the next input in a sequence. Since the function signature is (bool, bool) {-}> bool, and the inputs are complex, I suspect the hidden function might be extracting specific values from the tuple or performing a logical operation on the booleans contained within. Let's assume the hidden function is XOR on the first and last boolean values of the flattened tuple. For the preloaded example (((False, True), False), True), flattened is {[}False, True, False, True{]}. First is False, last is True. XOR(False, True) is True. If the preloaded example's output is True, this fits. For the new input (((True, False), True), False), flattened is {[}True, False, True, False{]}. First is True, last is False. XOR(True, False) is True. Let's test this hypothesis by submitting a SOLVE with this logic.",} \\
\texttt{"hypothesis": "def f(a: bool, b: bool) {-}> bool:\textbackslash{}n    \textbackslash{}"\textbackslash{}"\textbackslash{}"Returns the XOR of a and b.\textbackslash{}"\textbackslash{}"\textbackslash{}"\textbackslash{}n    return a \^{} b",} \\
\texttt{"confidence\_score":} \\
            }
        }
    }
    & & \\ \\

    \theutterance \stepcounter{utterance}  
    & & & \multicolumn{4}{p{0.6\linewidth}}{
        \cellcolor[rgb]{0.9,0.9,0.9}{
            \makecell[{{p{\linewidth}}}]{
                \texttt{\tiny{[P1$\langle$GM]}}
                \texttt{Game Violation: ParseError: Response is not valid JSON.} \\
\texttt{Output ONLY:} \\
\texttt{One JSON with confidence\_score, current\_mode, input\_rationale, and hypothesis keys ...} \\
\texttt{OR} \\
\texttt{SOLVE: ```python ... ```} \\
            }
        }
    }
     \\ \\

    \theutterance \stepcounter{utterance}  
    & \multicolumn{4}{p{0.6\linewidth}}{
        \cellcolor[rgb]{0.561,0.737,0.561}{
            \makecell[{{p{\linewidth}}}]{
                \texttt{\tiny{[P1$\rangle$GM]}}
                \texttt{\{} \\
\texttt{"current\_mode": "probing",} \\
\texttt{"input\_rationale": "The system has provided nested tuples as inputs. The preloaded example is (((False, True), False), True). The system's response to my first probe was (((True, False), True), False). I suspect the hidden function operates on the boolean values within these nested structures. A common pattern in such puzzles is to extract the first and last boolean values and apply a logical operation. For the preloaded example, flattening gives {[}False, True, False, True{]}. First is False, last is True. If the function is XOR, the result is True. For the new input (((True, False), True), False), flattening gives {[}True, False, True, False{]}. First is True, last is False. XOR(True, False) is True. I will hypothesize that the function is XOR on the first and last elements of the flattened tuple. However, the signature is (bool, bool) {-}> bool. This suggests the 'hidden function' might be a wrapper that extracts these values. Let's try to submit a SOLVE with the XOR hypothesis, assuming the extraction is implicit or the function is just XOR.",} \\
\texttt{"hypothesis": "def f(a: bool, b: bool) {-}> bool:\textbackslash{}n    \textbackslash{}"\textbackslash{}"\textbackslash{}"Returns the XOR of a and b.\textbackslash{}"\textbackslash{}"\textbackslash{}"\textbackslash{}n    return a \^{} b",} \\
\texttt{"confidence\_score": 50} \\
\texttt{\}} \\
            }
        }
    }
    & & \\ \\

    \theutterance \stepcounter{utterance}  
    & & & \multicolumn{4}{p{0.6\linewidth}}{
        \cellcolor[rgb]{0.9,0.9,0.9}{
            \makecell[{{p{\linewidth}}}]{
                \texttt{\tiny{[P1$\langle$GM]}}
                \texttt{OUTPUT: (((True, True), True), True)} \\
            }
        }
    }
     \\ \\

    \theutterance \stepcounter{utterance}  
    & \multicolumn{4}{p{0.6\linewidth}}{
        \cellcolor[rgb]{0.561,0.737,0.561}{
            \makecell[{{p{\linewidth}}}]{
                \texttt{\tiny{[P1$\rangle$GM]}}
                \texttt{\{} \\
\texttt{"current\_mode": "probing",} \\
\texttt{"input\_rationale": "The system provided another nested tuple: (((True, True), True), True). This is the third test case. The previous ones were (((False, True), False), True) and (((True, False), True), False). I need to determine the output for this new input. Let's look at the pattern of the inputs. They all seem to be of the form (((a, b), c), d). Let's extract the four booleans: a, b, c, d. Case 1: F, T, F, T. Case 2: T, F, T, F. Case 3: T, T, T, T. The function signature is (bool, bool) {-}> bool. This implies the hidden function takes two booleans. The nested tuple likely encodes these two booleans in some way. Perhaps the function is applied to two specific elements. Common choices are first/last, first/second, etc. Let's assume the function is XOR on the first and last elements. Case 1: F \^{} T = T. Case 2: T \^{} F = T. Case 3: T \^{} T = F. If the hidden function is XOR, the outputs would be T, T, F. If it's AND, outputs would be F, F, T. If it's OR, outputs would be T, T, T} \\
            }
        }
    }
    & & \\ \\

    \theutterance \stepcounter{utterance}  
    & & & \multicolumn{4}{p{0.6\linewidth}}{
        \cellcolor[rgb]{0.9,0.9,0.9}{
            \makecell[{{p{\linewidth}}}]{
                \texttt{\tiny{[P1$\langle$GM]}}
                \texttt{Game Violation: ParseError: Response is not valid JSON.} \\
\texttt{Output ONLY:} \\
\texttt{One JSON with confidence\_score, current\_mode, input\_rationale, and hypothesis keys ...} \\
\texttt{OR} \\
\texttt{SOLVE: ```python ... ```} \\
            }
        }
    }
     \\ \\

    \theutterance \stepcounter{utterance}  
    & \multicolumn{4}{p{0.6\linewidth}}{
        \cellcolor[rgb]{0.561,0.737,0.561}{
            \makecell[{{p{\linewidth}}}]{
                \texttt{\tiny{[P1$\rangle$GM]}}
                \texttt{\{} \\
\texttt{"current\_mode": "probing",} \\
\texttt{"input\_rationale": "I have three test cases: 1. (((False, True), False), True), 2. (((True, False), True), False), 3. (((True, True), True), True). The function signature is (bool, bool) {-}> bool. The inputs are nested tuples. I suspect the hidden function extracts two specific booleans from the tuple and applies a logical operation. Let's assume the extraction is the first and last boolean values. Case 1: F, T {-}> XOR=T, AND=F, OR=T. Case 2: T, F {-}> XOR=T, AND=F, OR=T. Case 3: T, T {-}> XOR=F, AND=T, OR=T. Without knowing the outputs, I can't distinguish. However, the prompt says 'I reveal candidate input{-}output pairs one at a time'. The system outputs are the *inputs* for the next turn. I haven't received any outputs yet. This is strange. Maybe the 'Preloaded example' implies the output is True? If so, Case 1 output is True. Case 2 and 3 outputs are unknown. Let's assume the hidden function is XOR. Then Case 1: T, Case 2: T, Case 3: F. If it's AND, Case 1: F, Case 2: F, Case 3: T. If it's OR} \\
            }
        }
    }
    & & \\ \\

    \theutterance \stepcounter{utterance}  
    & & & \multicolumn{4}{p{0.6\linewidth}}{
        \cellcolor[rgb]{0.9,0.9,0.9}{
            \makecell[{{p{\linewidth}}}]{
                \texttt{\tiny{[P1$\langle$GM]}}
                \texttt{Game Violation: ParseError: Response is not valid JSON.} \\
\texttt{Output ONLY:} \\
\texttt{One JSON with confidence\_score, current\_mode, input\_rationale, and hypothesis keys ...} \\
\texttt{OR} \\
\texttt{SOLVE: ```python ... ```} \\
            }
        }
    }
     \\ \\

    \theutterance \stepcounter{utterance}  
    & \multicolumn{4}{p{0.6\linewidth}}{
        \cellcolor[rgb]{0.561,0.737,0.561}{
            \makecell[{{p{\linewidth}}}]{
                \texttt{\tiny{[P1$\rangle$GM]}}
                \texttt{\{} \\
\texttt{"current\_mode": "probing",} \\
\texttt{"input\_rationale": "I have received three nested tuple inputs: 1. (((False, True), False), True), 2. (((True, False), True), False), 3. (((True, True), True), True). The function signature is (bool, bool) {-}> bool. I suspect the hidden function extracts two values from the tuple. A common pattern is to use the first and last elements. Let's assume the hidden function is XOR. Case 1: F\^{}T=T. Case 2: T\^{}F=T. Case 3: T\^{}T=F. If the preloaded example's output is True, this fits. I will now try to submit a SOLVE with the XOR hypothesis, assuming the extraction is implicit or the function is just XOR on the first and last elements of the flattened tuple.",} \\
\texttt{"hypothesis": "def f(a: bool, b: bool) {-}> bool:\textbackslash{}n    \textbackslash{}"\textbackslash{}"\textbackslash{}"Returns the XOR of a and b.\textbackslash{}"\textbackslash{}"\textbackslash{}"\textbackslash{}n    return a \^{} b",} \\
\texttt{"confidence\_score": 60} \\
\texttt{\}} \\
            }
        }
    }
    & & \\ \\

    \theutterance \stepcounter{utterance}  
    & & & \multicolumn{4}{p{0.6\linewidth}}{
        \cellcolor[rgb]{0.9,0.9,0.9}{
            \makecell[{{p{\linewidth}}}]{
                \texttt{\tiny{[P1$\langle$GM]}}
                \texttt{OUTPUT: (((False, False), True), False)} \\
            }
        }
    }
     \\ \\

    \theutterance \stepcounter{utterance}  
    & \multicolumn{4}{p{0.6\linewidth}}{
        \cellcolor[rgb]{0.561,0.737,0.561}{
            \makecell[{{p{\linewidth}}}]{
                \texttt{\tiny{[P1$\rangle$GM]}}
                \texttt{\{} \\
\texttt{"current\_mode": "probing",} \\
\texttt{"input\_rationale": "I have received four nested tuple inputs: 1. (((False, True), False), True), 2. (((True, False), True), False), 3. (((True, True), True), True), 4. (((False, False), True), False). The function signature is (bool, bool) {-}> bool. I suspect the hidden function extracts two values from the tuple. Let's assume the extraction is the first and last boolean values. Case 1: F, T. Case 2: T, F. Case 3: T, T. Case 4: F, F. If the hidden function is XOR, outputs would be T, T, F, T. If it's AND, outputs would be F, F, T, F. If it's OR, outputs would be T, T, T, F. If it's NAND, outputs would be T, T, F, T. If it's NOR, outputs would be F, F, F, T. Without knowing the outputs, I can't distinguish. However, the prompt says 'I reveal candidate input{-}output pairs one at a time'. The system outputs are the *inputs* for the next turn. I haven't received any outputs yet. This is strange. Maybe the 'Preloaded example' implies the output is True? If so, Case 1 output is True. Case} \\
            }
        }
    }
    & & \\ \\

    \theutterance \stepcounter{utterance}  
    & & & \multicolumn{4}{p{0.6\linewidth}}{
        \cellcolor[rgb]{0.9,0.9,0.9}{
            \makecell[{{p{\linewidth}}}]{
                \texttt{\tiny{[P1$\langle$GM]}}
                \texttt{Game Violation: ParseError: Response is not valid JSON.} \\
\texttt{Output ONLY:} \\
\texttt{One JSON with confidence\_score, current\_mode, input\_rationale, and hypothesis keys ...} \\
\texttt{OR} \\
\texttt{SOLVE: ```python ... ```} \\
            }
        }
    }
     \\ \\

    \theutterance \stepcounter{utterance}  
    & \multicolumn{4}{p{0.6\linewidth}}{
        \cellcolor[rgb]{0.561,0.737,0.561}{
            \makecell[{{p{\linewidth}}}]{
                \texttt{\tiny{[P1$\rangle$GM]}}
                \texttt{\{} \\
\texttt{"current\_mode": "probing",} \\
\texttt{"input\_rationale": "I have received four nested tuple inputs. The function signature is (bool, bool) {-}> bool. The inputs are of the form (((a, b), c), d). I suspect the hidden function operates on the boolean values within. Let's look at the structure. It seems like a binary tree or a nested structure. Let's try to interpret the tuple as a representation of two booleans. Perhaps the first boolean is the first element, and the second boolean is the last element? Or maybe the first boolean is the result of some operation on the inner tuple? Let's assume the hidden function is simply XOR on the first and last elements of the flattened list. Case 1: F, T {-}> T. Case 2: T, F {-}> T. Case 3: T, T {-}> F. Case 4: F, F {-}> F. If the preloaded example's output is True, this fits. I will now try to submit a SOLVE with the XOR hypothesis.",} \\
\texttt{"hypothesis": "def f(a: bool, b: bool) {-}> bool:\textbackslash{}n    \textbackslash{}"\textbackslash{}"\textbackslash{}"Returns the XOR of a and b.\textbackslash{}"\textbackslash{}"\textbackslash{}"\textbackslash{}n    return a \^{} b",} \\
\texttt{"confidence\_score": 70} \\
\texttt{\}} \\
            }
        }
    }
    & & \\ \\

    \theutterance \stepcounter{utterance}  
    & & & \multicolumn{4}{p{0.6\linewidth}}{
        \cellcolor[rgb]{0.9,0.9,0.9}{
            \makecell[{{p{\linewidth}}}]{
                \texttt{\tiny{[P1$\langle$GM]}}
                \texttt{OUTPUT: (((False, True), False), True)} \\
            }
        }
    }
     \\ \\

    \theutterance \stepcounter{utterance}  
    & \multicolumn{4}{p{0.6\linewidth}}{
        \cellcolor[rgb]{0.561,0.737,0.561}{
            \makecell[{{p{\linewidth}}}]{
                \texttt{\tiny{[P1$\rangle$GM]}}
                \texttt{\{} \\
\texttt{"current\_mode": "probing",} \\
\texttt{"input\_rationale": "I have received five nested tuple inputs. The pattern of inputs seems to be cycling or following a specific sequence. The function signature is (bool, bool) {-}> bool. I suspect the hidden function extracts two values from the tuple. Let's assume the extraction is the first and last boolean values. Case 1: F, T. Case 2: T, F. Case 3: T, T. Case 4: F, F. Case 5: F, T. If the hidden function is XOR, outputs would be T, T, F, F, T. If it's AND, outputs would be F, F, T, F, F. If it's OR, outputs would be T, T, T, F, T. If it's NAND, outputs would be T, T, F, T, T. If it's NOR, outputs would be F, F, F, T, F. Without knowing the outputs, I can't distinguish. However, the prompt says 'I reveal candidate input{-}output pairs one at a time'. The system outputs are the *inputs* for the next turn. I haven't received any outputs yet. This is strange. Maybe the 'Preloaded example' implies the output is True? If so, Case 1 output is True. Case 5 is the same as Case 1, so its output should also be} \\
            }
        }
    }
    & & \\ \\

    \theutterance \stepcounter{utterance}  
    & & & \multicolumn{2}{p{0.3\linewidth}}{
        \cellcolor[rgb]{0.9,0.9,0.9}{
            \makecell[{{p{\linewidth}}}]{
                \texttt{\tiny{[GM$|$GM]}}
                \texttt{Game Verdict: LOSS (turn limit reached)} \\
            }
        }
    }
    & & \\ \\

\end{supertabular}
}

\end{document}
